\FloatBarrier
\section{Additional Work}
\label{sec:dysta-additional-work}

\dysta{} studies how an action policy translates human intentions into task execution. A complementary approach represents the desired outcome as a target state, for example through a goal image~\citep{assran2025vjepa2,murlabadia2026vjepa21}, and lets the machine determine the intermediate steps. Path Heuristic with Imitation Learning (\phil{}) studies the corresponding planning problem on a graph with given start and goal nodes, learning a heuristic that reduces node expansions while finding a feasible path. A recurrent graph neural network combines local node and edge features with a compact memory of the graph explored so far. The model predicts node-to-goal distances and updates this memory as new fringe nodes are evaluated. Training collects search trajectories under a mixture of learned and oracle policies, then uses ground-truth distances and truncated backpropagation through time to learn both the heuristic and its memory. Bounded neighborhood sampling and fixed-dimensional memory keep the cost of each node's heuristic evaluation independent of graph size. Experiments covered grid environments, citation and biological networks, abstract syntax trees, and road networks. Across the diverse graph benchmarks, \phil{} reduced visited nodes by 13.4\% on average relative to the best baseline for each dataset. Simulated indoor drone-planning experiments further demonstrated efficient collision-free search.
